\chapter{Pixels, shuffle controls and scalar tails}
\MARSChapterHeading{Pixels, shuffle controls and tails}
\label{revision-lab:msa}

Use \texttt{rgb-formats.asm} in \texttt{revision/labs/files} for observation and
\texttt{rgb555.asm} for independent writing. Read the main chapter's lane model
and RGB565 field derivation first. These activities use ordinary Run and Debug;
a native guided version of this unit remains unfinished. Keep disposable copies
for intentional faults so that each comparison has one known change.

\section{Separate a pixel value from its display adapter}

The first eight RGB triples in the worked source are deliberately varied. Derive
the packed words for $(123,45,67)$, black, white and $(128,64,247)$. Predict their
stored byte pairs before running. At \texttt{pack\_ready}, check the returned count
and inspect \texttt{packed\_a}. The caller has not yet stored
\texttt{packed\_count}; distinguish the function's result from its later publication.

Continue to \texttt{unpack\_ready}. The first expanded triple is $(123,44,66)$.
Explain which low bits were lost and how repetition fills the expanded fields.
Test red 128 separately: the expanded value is 132. Packing and expansion are
specified code-value operations; neither an exact inverse on RGB888 nor an arbitrary
interpolation gives the required result.

At \texttt{preview\_ready}, the top row uses original pixels and the bottom row
expanded pixels. The preview holds \texttt{00RRGGBB} display words. It is a separate
adapter: the Bitmap Display does not reinterpret raw RGB565 bytes as those words.
A small visual difference cannot substitute for inspecting the numerical result.

\section{Predict the first gathered vector}

Break at \texttt{pack\_loaded}, after the two loads but before the control copies.
Expand MSA vector registers and inspect \texttt{\$w0} and \texttt{\$w1} as byte
lanes. The first begins with \texttt{7b 2d 43}; the second begins with the byte at
source offset eight, not sixteen. Derive all 24 source byte positions covered by
the two windows. Mark the overlap before interpreting shuffle selectors.

Continue to \texttt{pack\_gathered}. Expand \texttt{\$w4} and select its
sixteen-bit integer lane row. In lane order, the gathered red values are
$[123,0,255,128,7,8,248,31]$. The same bits in byte lanes alternate components
and zero bytes. These are two views of one register, not two stored copies.
Derive the green and blue halfwords and check them in \texttt{\$w5/\$w6}.
The debugger's integer lane rows use hexadecimal; convert the expected values
before deciding that a displayed value is wrong.

Locate red pixel six at input offset eighteen. Explain why its control is 26.
Locate green pixel five at offset sixteen and explain why its control is 24.
Changing only the source operands' order in one shuffle should change those
selections. Investigate the first incorrect gathered lane before continuing to
packed output. A field-shift change cannot repair a wrong component source.

\section{Follow field placement and preserve controls}

At \texttt{pack\_fields}, the shifts and ORs have completed but the vector store
has not. The first eight halfwords are
\texttt{7968, 0000, ffff, 821e, 0000, 0821, ffff, 1bfb}.
Trace lane zero from 123, 45 and 67 through its three fields. Then inspect the
first store and compare its sixteen bytes with these eight halfwords. All
sixteen output bytes belong to this group, even though its input occupied 24.

\begin{figure}[htbp]
\centering
\begin{tikzpicture}[alt={Two actual light-theme debugger views of vector register w4. Before the shifts, its halfword row is 007b, 0000, 00ff, 0080, 0007, 0008, 00f8, 001f: eight gathered red components. Before the store, it is 7968, 0000, ffff, 821e, 0000, 0821, ffff, 1bfb: eight packed RGB565 values. Each capture also shows the same register as sixteen byte lanes. Native boxes identify the halfword rows.}]
\node[anchor=north west,inner sep=0,font=\small\bfseries] (atitle) at (0,0) {1. At \texttt{pack\_gathered}, before the shifts};
\node[anchor=north west,inner sep=0] (a) at ([yshift=-2mm]atitle.south west) {\MARSRevisionCrop{revision/captures/vector/rgb-gathered.png}{0.97\linewidth}{0.23125}{0.180}{0.2625}{0.758}};
\node[anchor=north west,inner sep=0,font=\small\bfseries] (btitle) at ([yshift=-5mm]a.south west) {2. At \texttt{pack\_fields}, before the store};
\node[anchor=north west,inner sep=0] (b) at ([yshift=-2mm]btitle.south west) {\MARSRevisionCrop{revision/captures/vector/rgb-fields.png}{0.97\linewidth}{0.23125}{0.180}{0.2625}{0.758}};
\foreach \panel/\rightedge in {a/.74,b/.715} {
 \begin{scope}[shift={(\panel.south west)},x={($(\panel.south east)-(\panel.south west)$)},y={($(\panel.north west)-(\panel.south west)$)}]
 \draw[black,line width=.7pt] (.038,.025) rectangle (\rightedge,.36);
 \end{scope}
}
\end{tikzpicture}
\caption{Actual native Debug observations of the same first group. Expand MSA vector registers, then \texttt{\$w4}. Its byte and halfword rows reinterpret the same bits. Native boxes mark the halfword row before arithmetic and before output. In the second view, the first byte pair \texttt{68 79} and halfword \texttt{7968} describe the same packed pixel.}
\label{fig:revision-msa-lanes}
\end{figure}


Use at least sixteen pixels to observe two group iterations. Record
\texttt{\$w16} before and after the first group: it retains the red controls.
Record \texttt{\$w4}: it contains results. In a copy, skip the control copies on
the second iteration and inspect the changed gather. Restore the original before
checking other faults. Passing the first eight pixels does not validate control
reuse across a longer input.

Do not assume that every different suffix creates a fault. Consider replacing
the red left shift with a word-lane shift by eleven. Before that shift, each red
halfword is at most 31; its placed bits still fit its original sixteen-bit lane.
Predict whether any surviving bit crosses the boundary, then check the result.
In contrast, putting an input component into a green halfword's high byte can
leave unwanted bits after the green shifts. Trace which bits survive before
judging equivalence.

\section{Prove the gate and the remainder together}

Use counts zero, seven, eight, nine and seventeen. For each, predict the number
of full groups and scalar pixels. Inspect \texttt{\$a0}, \texttt{\$a1} and
\texttt{\$a2} at each \texttt{pack\_block} visit and at \texttt{pack\_tail}.
The source advances by 24 bytes per group and output by sixteen; \texttt{\$v0}
keeps the original count. At count zero there are no caller-data accesses.

Move the second load from offset eight to sixteen in a copy. For an eight-pixel
input, identify the first byte beyond the selected source. That byte can be mapped
and still violate the contract. Restore the load, then change the eligibility
test to allow seven pixels. Identify the load whose full range no longer fits.
The public checks inspect access footprints, not just final output bytes.

Run the original public interface and caller checks against the worked source:
\begin{verbatim}
java -jar /path/to/cli.jar grade \
  packages/019-rgb-formats/checks/kernel.kts \
  revision/labs/files/rgb-formats.asm
java -jar /path/to/cli.jar grade \
  packages/019-rgb-formats/checks/program.kts \
  revision/labs/files/rgb-formats.asm
\end{verbatim}
These checks accept equivalent scalar and vector implementations. They do not
establish an MSA performance advantage or require the worked source's register
assignments. If you compare costs, state the interval and complete workload as in
Chapter~\ref{revision-lab:measurement}; a vector instruction can move more bytes
than its scalar counterpart.

\section{Implement RGB555 without a construction sequence}

Open \texttt{rgb555.asm} and implement \texttt{pack\_rgb555}. Keep its supplied
caller and data harness. Arguments are full input and output pointers plus count
in \texttt{\$a0/\$a1/\$a2}; return count in \texttt{\$v0}. Count is zero through
32; input has $3n$ bytes and output $2n$. They are disjoint, may begin at any byte
alignment and may select interior windows. Read only selected caller input or
output, and write only selected output. Preserve inputs, metadata, guards,
inactive output, the full saved registers, frame pointer and stack pointer.
Private constants and native stack storage are allowed. Return through the
supplied continuation and perform no console I/O.

Derive the five-bit positions from the main chapter's equation. The top bit is
zero. The first default pixel becomes \texttt{3ca8}, stored \texttt{a8 3c}; white
becomes \texttt{7fff}. Choose a further green value that distinguishes RGB555
from RGB565. Establish scalar correctness before choosing a vector group and tail.
Explain whether the existing channel gather needs to change when only the output
field positions change.

Run the public checks:
\begin{verbatim}
java -jar /path/to/cli.jar grade \
  revision/labs/checks/rgb555.kts \
  /path/to/your/rgb555.asm
\end{verbatim}
The suite checks every supported count with varied byte offsets and a default
caller run. It accepts equivalent implementations. Keep a scalar result, a
vector-group observation, an incomplete-group trace and your bounds argument.
If an instructor assigns vectorization, that explanation and observation establish
the strategy separately from behavioral correctness. A successful output comparison
does not show which instructions you chose.
